\documentclass[10pt]{letter}
\usepackage[margin=0.8in]{geometry}
\usepackage{parskip}
\signature{Ranyang Li (on behalf of all authors)}
\address{School of Artificial Intelligence and Big Data\\Henan University of Technology\\Zhengzhou 450001, China\\lry@haut.edu.cn}
\begin{document}
\begin{letter}{The Editor-in-Chief\\IEEE Transactions on Medical Imaging}
\opening{Dear Editor,}

We have prepared the following original research manuscript for consideration
for publication in \emph{IEEE Transactions on Medical Imaging}:

\textbf{``FreeReg: Free-Coordinate Registration of 3D Models for Augmented Reality in Image-Guided Bronchoscopy''}

Image-guided bronchoscopy and transbronchial biopsy depend on registering the
preoperative CT airway model onto the live endoscopic view. Deployed
electromagnetic navigation bronchoscopy still reports 5--10\,mm localization
error, loses track under respiration and coughing, and requires an external
field generator and tracked probe. Classical calibrate-then-align registration
assumes a calibrated camera, a rigid target, and independent per-frame
estimation; all three fail in the airway. We formulate CT-to-bronchoscopy
registration as free-coordinate AR/MR registration: an explicit, evidence-
supported mapping from the native 3-D model frame to camera and pixel
coordinates, rather than unconstrained visual overlay. We solve it with a
centerline-driven correspondence engine, an observability-guided
rigid-plus-warp estimator that deforms only within the provably observable
subspace, a manifold SE(3) Kalman filter with innovation gating, and a
co-anchored registration--rendering field whose photometric residual is an
online validation signal.

Key results: (1)~on 12 clinical CT airways (270 synthetic fly-through
frames with ground-truth poses) the pipeline registers at 4.23\,mm mean
target registration error (92\% $<$10\,mm, Wilcoxon $p=4.4\times10^{-31}$)
at 30--80\,FPS, within the 5--10\,mm error range of electromagnetic
navigation bronchoscopy and without external tracking hardware; real
intra-operative video is used for qualitative AR demonstration;
(2)~observability-guided non-rigid estimation recovers deformation to 1.1\% on
a real-colon FEM benchmark and provides, to our knowledge, the first explicit
account of which airway deformation directions monocular endoscopy can and
cannot recover (rank deficiency 151/648 directions, 23\%); (3)~as a
correctness/generalization check, the same engine matches the pose champion
FoundationPose on the BOP Linemod benchmark (both at ADD-S 1.000 on the full
2600-sample benchmark) while remaining the only method that stays accurate in
the airway, where FoundationPose (TRE 104\,mm) and Endo3R (ATE 34\,mm) collapse
out-of-distribution; (4)~a learned front-end refined by ICP matches the
geometric engine at $5\times$ speed (ADD-S 0.987 at 158\,ms).

The work aligns with IEEE TMI's scope in medical image registration,
image-guided intervention, and computational methodology for minimally invasive
diagnosis. This manuscript is original, has not been published previously, and
is not under consideration elsewhere. All authors have approved the manuscript
and declare no competing interests. The clinical data statement is provided in
the accompanying Ethics Statement; the study is a retrospective secondary
analysis of de-identified clinical imaging authorised through the institutional
data-governance process of Henan Provincial People's Hospital.

Thank you for your consideration.

\closing{Sincerely,}
\end{letter}
\end{document}
